\input eplain
\parindent=0pt
\pdfpagewidth=148mm
\pdfpageheight=210mm
\hsize=128mm
\vsize=190mm
\pdfhorigin=0pt
\pdfvorigin=0pt
\hoffset=10mm
\voffset=10mm
\xrdef{g0}
/
\par\begingroup
\leftskip=2em
\smallskip
calculus \leaders\hbox to 10pt {\hss.\hss}  \hfill \xref{g1}
\smallskip
ink \leaders\hbox to 10pt {\hss.\hss}  \hfill \xref{g2}
\smallskip
\medskip
\par\endgroup
\bigskip
\xrdef{g1}
calculus
\par\begingroup
\leftskip=2em
\smallskip
pearson \leaders\hbox to 10pt {\hss.\hss}  \hfill \xref{g3}
\smallskip
\medskip
\par\endgroup
\bigskip
\xrdef{g2}
ink
\par\begingroup
\leftskip=2em
\smallskip
\medskip
07C \leaders\hbox to 10pt {\hss.\hss}  \hfill \xref{r59}
\smallskip
07D \leaders\hbox to 10pt {\hss.\hss}  \hfill \xref{r60}
\smallskip
07I \leaders\hbox to 10pt {\hss.\hss}  \hfill \xref{r61}
\smallskip
07M \leaders\hbox to 10pt {\hss.\hss}  \hfill \xref{r62}
\smallskip
07O \leaders\hbox to 10pt {\hss.\hss}  \hfill \xref{r63}
\smallskip
07V \leaders\hbox to 10pt {\hss.\hss}  \hfill \xref{r64}
\smallskip
07W \leaders\hbox to 10pt {\hss.\hss}  \hfill \xref{r65}
\smallskip
07X \leaders\hbox to 10pt {\hss.\hss}  \hfill \xref{r66}
\smallskip
07Y \leaders\hbox to 10pt {\hss.\hss}  \hfill \xref{r67}
\smallskip
080 \leaders\hbox to 10pt {\hss.\hss}  \hfill \xref{r68}
\smallskip
\par\endgroup
\bigskip
\xrdef{g3}
calculus/pearson
\par\begingroup
\leftskip=2em
\smallskip
examples \leaders\hbox to 10pt {\hss.\hss}  \hfill \xref{g14}
\smallskip
theorems \leaders\hbox to 10pt {\hss.\hss}  \hfill \xref{g17}
\smallskip
toc \leaders\hbox to 10pt {\hss.\hss}  \hfill \xref{g18}
\smallskip
\medskip
ink \leaders\hbox to 10pt {\hss.\hss}  \hfill \xref{r13}
\smallskip
news \leaders\hbox to 10pt {\hss.\hss}  \hfill \xref{r14}
\smallskip
theorems \leaders\hbox to 10pt {\hss.\hss}  \hfill \xref{r15}
\smallskip
toc \leaders\hbox to 10pt {\hss.\hss}  \hfill \xref{r49}
\smallskip
\par\endgroup
\bigskip
\xrdef{g14}
calculus/pearson/examples
\par\begingroup
\leftskip=2em
\smallskip
\medskip
2 4 01 \leaders\hbox to 10pt {\hss.\hss}  \hfill \xref{r1}
\smallskip
2 4 02 \leaders\hbox to 10pt {\hss.\hss}  \hfill \xref{r2}
\smallskip
2 4 03 \leaders\hbox to 10pt {\hss.\hss}  \hfill \xref{r3}
\smallskip
2 4 04 \leaders\hbox to 10pt {\hss.\hss}  \hfill \xref{r4}
\smallskip
2 4 05 \leaders\hbox to 10pt {\hss.\hss}  \hfill \xref{r5}
\smallskip
2 4 06 \leaders\hbox to 10pt {\hss.\hss}  \hfill \xref{r6}
\smallskip
2 5 01 \leaders\hbox to 10pt {\hss.\hss}  \hfill \xref{r7}
\smallskip
2 5 02 \leaders\hbox to 10pt {\hss.\hss}  \hfill \xref{r8}
\smallskip
2 5 03 \leaders\hbox to 10pt {\hss.\hss}  \hfill \xref{r9}
\smallskip
2 5 04 \leaders\hbox to 10pt {\hss.\hss}  \hfill \xref{r10}
\smallskip
2 5 05 \leaders\hbox to 10pt {\hss.\hss}  \hfill \xref{r11}
\smallskip
2 5 06 \leaders\hbox to 10pt {\hss.\hss}  \hfill \xref{r12}
\smallskip
\par\endgroup
\bigskip
\xrdef{g17}
calculus/pearson/theorems
\par\begingroup
\leftskip=2em
\smallskip
\medskip
2 10 polynomial and rational functions \leaders\hbox to 10pt {\hss.\hss}  \hfill \xref{r16}
\smallskip
2 11 continunity of composite functions \leaders\hbox to 10pt {\hss.\hss}  \hfill \xref{r17}
\smallskip
2 12 limits of composite functions \leaders\hbox to 10pt {\hss.\hss}  \hfill \xref{r18}
\smallskip
2 13 continuity of functions with roots \leaders\hbox to 10pt {\hss.\hss}  \hfill \xref{r19}
\smallskip
2 14 continuity of inverse functions \leaders\hbox to 10pt {\hss.\hss}  \hfill \xref{r20}
\smallskip
2 15 continuity of transcendental \leaders\hbox to 10pt {\hss.\hss}  \hfill \xref{r21}
\smallskip
2 16 intermediate value theorem \leaders\hbox to 10pt {\hss.\hss}  \hfill \xref{r22}
\smallskip
2 1 relationship between one-sided \leaders\hbox to 10pt {\hss.\hss}  \hfill \xref{r23}
\smallskip
2 2 limits of a linear function \leaders\hbox to 10pt {\hss.\hss}  \hfill \xref{r24}
\smallskip
2 3a limit law: sum \leaders\hbox to 10pt {\hss.\hss}  \hfill \xref{r25}
\smallskip
2 3b limit law: difference \leaders\hbox to 10pt {\hss.\hss}  \hfill \xref{r26}
\smallskip
2 3c limit law: constant multiple \leaders\hbox to 10pt {\hss.\hss}  \hfill \xref{r27}
\smallskip
2 3d limit law: product \leaders\hbox to 10pt {\hss.\hss}  \hfill \xref{r28}
\smallskip
2 3e limit law: quotient \leaders\hbox to 10pt {\hss.\hss}  \hfill \xref{r29}
\smallskip
2 3f limit law: power \leaders\hbox to 10pt {\hss.\hss}  \hfill \xref{r30}
\smallskip
2 3g limit law: root \leaders\hbox to 10pt {\hss.\hss}  \hfill \xref{r31}
\smallskip
2 3h limit laws for one-sided limits \leaders\hbox to 10pt {\hss.\hss}  \hfill \xref{r32}
\smallskip
2 4 limits of polynomial and rational \leaders\hbox to 10pt {\hss.\hss}  \hfill \xref{r33}
\smallskip
2 5 the squeeze theorem \leaders\hbox to 10pt {\hss.\hss}  \hfill \xref{r34}
\smallskip
2 6 limits at infinity of powers and \leaders\hbox to 10pt {\hss.\hss}  \hfill \xref{r35}
\smallskip
2 7 end behavior and asymptotes of \leaders\hbox to 10pt {\hss.\hss}  \hfill \xref{r36}
\smallskip
2 8 end behavior of e x, e {-x}, and \leaders\hbox to 10pt {\hss.\hss}  \hfill \xref{r37}
\smallskip
2 9 continuity rules \leaders\hbox to 10pt {\hss.\hss}  \hfill \xref{r38}
\smallskip
3 1 differentiable implies continuous \leaders\hbox to 10pt {\hss.\hss}  \hfill \xref{r39}
\smallskip
3 1a not continuous implies not differentiable \leaders\hbox to 10pt {\hss.\hss}  \hfill \xref{r40}
\smallskip
3 2 constant rule \leaders\hbox to 10pt {\hss.\hss}  \hfill \xref{r41}
\smallskip
3 3 power rule \leaders\hbox to 10pt {\hss.\hss}  \hfill \xref{r42}
\smallskip
3 4 constant multiple rule \leaders\hbox to 10pt {\hss.\hss}  \hfill \xref{r43}
\smallskip
3 5 sum rule \leaders\hbox to 10pt {\hss.\hss}  \hfill \xref{r44}
\smallskip
3 6 the derivative of e x \leaders\hbox to 10pt {\hss.\hss}  \hfill \xref{r45}
\smallskip
3 7 product rule \leaders\hbox to 10pt {\hss.\hss}  \hfill \xref{r46}
\smallskip
3 8 quotient rule \leaders\hbox to 10pt {\hss.\hss}  \hfill \xref{r47}
\smallskip
3 9 power rule \leaders\hbox to 10pt {\hss.\hss}  \hfill \xref{r48}
\smallskip
\par\endgroup
\bigskip
\xrdef{g18}
calculus/pearson/toc
\par\begingroup
\leftskip=2em
\smallskip
02 limits \leaders\hbox to 10pt {\hss.\hss}  \hfill \xref{g64}
\smallskip
03 derivatives \leaders\hbox to 10pt {\hss.\hss}  \hfill \xref{g65}
\smallskip
\medskip
02 limits \leaders\hbox to 10pt {\hss.\hss}  \hfill \xref{r50}
\smallskip
03 derivatives \leaders\hbox to 10pt {\hss.\hss}  \hfill \xref{r54}
\smallskip
\par\endgroup
\bigskip
\xrdef{g64}
calculus/pearson/toc/02 limits
\par\begingroup
\leftskip=2em
\smallskip
\medskip
04 infinite limits \leaders\hbox to 10pt {\hss.\hss}  \hfill \xref{r51}
\smallskip
05 limits at infinity \leaders\hbox to 10pt {\hss.\hss}  \hfill \xref{r52}
\smallskip
06 continuity \leaders\hbox to 10pt {\hss.\hss}  \hfill \xref{r53}
\smallskip
\par\endgroup
\bigskip
\xrdef{g65}
calculus/pearson/toc/03 derivatives
\par\begingroup
\leftskip=2em
\smallskip
\medskip
01 introducing the derivative \leaders\hbox to 10pt {\hss.\hss}  \hfill \xref{r55}
\smallskip
02 derivative as function \leaders\hbox to 10pt {\hss.\hss}  \hfill \xref{r56}
\smallskip
03 rules of differentiation \leaders\hbox to 10pt {\hss.\hss}  \hfill \xref{r57}
\smallskip
04 product and quotient rules \leaders\hbox to 10pt {\hss.\hss}  \hfill \xref{r58}
\smallskip
\par\endgroup
\bigskip
calculus/pearson/examples/2 4 01 \leaders\hbox to 10pt {\hss.\hss}  \hfill \xref{r1}
\smallskip
calculus/pearson/examples/2 4 02 \leaders\hbox to 10pt {\hss.\hss}  \hfill \xref{r2}
\smallskip
calculus/pearson/examples/2 4 03 \leaders\hbox to 10pt {\hss.\hss}  \hfill \xref{r3}
\smallskip
calculus/pearson/examples/2 4 04 \leaders\hbox to 10pt {\hss.\hss}  \hfill \xref{r4}
\smallskip
calculus/pearson/examples/2 4 05 \leaders\hbox to 10pt {\hss.\hss}  \hfill \xref{r5}
\smallskip
calculus/pearson/examples/2 4 06 \leaders\hbox to 10pt {\hss.\hss}  \hfill \xref{r6}
\smallskip
calculus/pearson/examples/2 5 01 \leaders\hbox to 10pt {\hss.\hss}  \hfill \xref{r7}
\smallskip
calculus/pearson/examples/2 5 02 \leaders\hbox to 10pt {\hss.\hss}  \hfill \xref{r8}
\smallskip
calculus/pearson/examples/2 5 03 \leaders\hbox to 10pt {\hss.\hss}  \hfill \xref{r9}
\smallskip
calculus/pearson/examples/2 5 04 \leaders\hbox to 10pt {\hss.\hss}  \hfill \xref{r10}
\smallskip
calculus/pearson/examples/2 5 05 \leaders\hbox to 10pt {\hss.\hss}  \hfill \xref{r11}
\smallskip
calculus/pearson/examples/2 5 06 \leaders\hbox to 10pt {\hss.\hss}  \hfill \xref{r12}
\smallskip
calculus/pearson/ink \leaders\hbox to 10pt {\hss.\hss}  \hfill \xref{r13}
\smallskip
calculus/pearson/news \leaders\hbox to 10pt {\hss.\hss}  \hfill \xref{r14}
\smallskip
calculus/pearson/theorems \leaders\hbox to 10pt {\hss.\hss}  \hfill \xref{r15}
\smallskip
calculus/pearson/theorems/2 10 polynomial and rational functions \leaders\hbox to 10pt {\hss.\hss}  \hfill \xref{r16}
\smallskip
calculus/pearson/theorems/2 11 continunity of composite functions \leaders\hbox to 10pt {\hss.\hss}  \hfill \xref{r17}
\smallskip
calculus/pearson/theorems/2 12 limits of composite functions \leaders\hbox to 10pt {\hss.\hss}  \hfill \xref{r18}
\smallskip
calculus/pearson/theorems/2 13 continuity of functions with roots \leaders\hbox to 10pt {\hss.\hss}  \hfill \xref{r19}
\smallskip
calculus/pearson/theorems/2 14 continuity of inverse functions \leaders\hbox to 10pt {\hss.\hss}  \hfill \xref{r20}
\smallskip
calculus/pearson/theorems/2 15 continuity of transcendental \leaders\hbox to 10pt {\hss.\hss}  \hfill \xref{r21}
\smallskip
calculus/pearson/theorems/2 16 intermediate value theorem \leaders\hbox to 10pt {\hss.\hss}  \hfill \xref{r22}
\smallskip
calculus/pearson/theorems/2 1 relationship between one-sided \leaders\hbox to 10pt {\hss.\hss}  \hfill \xref{r23}
\smallskip
calculus/pearson/theorems/2 2 limits of a linear function \leaders\hbox to 10pt {\hss.\hss}  \hfill \xref{r24}
\smallskip
calculus/pearson/theorems/2 3a limit law: sum \leaders\hbox to 10pt {\hss.\hss}  \hfill \xref{r25}
\smallskip
calculus/pearson/theorems/2 3b limit law: difference \leaders\hbox to 10pt {\hss.\hss}  \hfill \xref{r26}
\smallskip
calculus/pearson/theorems/2 3c limit law: constant multiple \leaders\hbox to 10pt {\hss.\hss}  \hfill \xref{r27}
\smallskip
calculus/pearson/theorems/2 3d limit law: product \leaders\hbox to 10pt {\hss.\hss}  \hfill \xref{r28}
\smallskip
calculus/pearson/theorems/2 3e limit law: quotient \leaders\hbox to 10pt {\hss.\hss}  \hfill \xref{r29}
\smallskip
calculus/pearson/theorems/2 3f limit law: power \leaders\hbox to 10pt {\hss.\hss}  \hfill \xref{r30}
\smallskip
calculus/pearson/theorems/2 3g limit law: root \leaders\hbox to 10pt {\hss.\hss}  \hfill \xref{r31}
\smallskip
calculus/pearson/theorems/2 3h limit laws for one-sided limits \leaders\hbox to 10pt {\hss.\hss}  \hfill \xref{r32}
\smallskip
calculus/pearson/theorems/2 4 limits of polynomial and rational \leaders\hbox to 10pt {\hss.\hss}  \hfill \xref{r33}
\smallskip
calculus/pearson/theorems/2 5 the squeeze theorem \leaders\hbox to 10pt {\hss.\hss}  \hfill \xref{r34}
\smallskip
calculus/pearson/theorems/2 6 limits at infinity of powers and \leaders\hbox to 10pt {\hss.\hss}  \hfill \xref{r35}
\smallskip
calculus/pearson/theorems/2 7 end behavior and asymptotes of \leaders\hbox to 10pt {\hss.\hss}  \hfill \xref{r36}
\smallskip
calculus/pearson/theorems/2 8 end behavior of e x, e {-x}, and \leaders\hbox to 10pt {\hss.\hss}  \hfill \xref{r37}
\smallskip
calculus/pearson/theorems/2 9 continuity rules \leaders\hbox to 10pt {\hss.\hss}  \hfill \xref{r38}
\smallskip
calculus/pearson/theorems/3 1 differentiable implies continuous \leaders\hbox to 10pt {\hss.\hss}  \hfill \xref{r39}
\smallskip
calculus/pearson/theorems/3 1a not continuous implies not differentiable \leaders\hbox to 10pt {\hss.\hss}  \hfill \xref{r40}
\smallskip
calculus/pearson/theorems/3 2 constant rule \leaders\hbox to 10pt {\hss.\hss}  \hfill \xref{r41}
\smallskip
calculus/pearson/theorems/3 3 power rule \leaders\hbox to 10pt {\hss.\hss}  \hfill \xref{r42}
\smallskip
calculus/pearson/theorems/3 4 constant multiple rule \leaders\hbox to 10pt {\hss.\hss}  \hfill \xref{r43}
\smallskip
calculus/pearson/theorems/3 5 sum rule \leaders\hbox to 10pt {\hss.\hss}  \hfill \xref{r44}
\smallskip
calculus/pearson/theorems/3 6 the derivative of e x \leaders\hbox to 10pt {\hss.\hss}  \hfill \xref{r45}
\smallskip
calculus/pearson/theorems/3 7 product rule \leaders\hbox to 10pt {\hss.\hss}  \hfill \xref{r46}
\smallskip
calculus/pearson/theorems/3 8 quotient rule \leaders\hbox to 10pt {\hss.\hss}  \hfill \xref{r47}
\smallskip
calculus/pearson/theorems/3 9 power rule \leaders\hbox to 10pt {\hss.\hss}  \hfill \xref{r48}
\smallskip
calculus/pearson/toc \leaders\hbox to 10pt {\hss.\hss}  \hfill \xref{r49}
\smallskip
calculus/pearson/toc/02 limits \leaders\hbox to 10pt {\hss.\hss}  \hfill \xref{r50}
\smallskip
calculus/pearson/toc/02 limits/04 infinite limits \leaders\hbox to 10pt {\hss.\hss}  \hfill \xref{r51}
\smallskip
calculus/pearson/toc/02 limits/05 limits at infinity \leaders\hbox to 10pt {\hss.\hss}  \hfill \xref{r52}
\smallskip
calculus/pearson/toc/02 limits/06 continuity \leaders\hbox to 10pt {\hss.\hss}  \hfill \xref{r53}
\smallskip
calculus/pearson/toc/03 derivatives \leaders\hbox to 10pt {\hss.\hss}  \hfill \xref{r54}
\smallskip
calculus/pearson/toc/03 derivatives/01 introducing the derivative \leaders\hbox to 10pt {\hss.\hss}  \hfill \xref{r55}
\smallskip
calculus/pearson/toc/03 derivatives/02 derivative as function \leaders\hbox to 10pt {\hss.\hss}  \hfill \xref{r56}
\smallskip
calculus/pearson/toc/03 derivatives/03 rules of differentiation \leaders\hbox to 10pt {\hss.\hss}  \hfill \xref{r57}
\smallskip
calculus/pearson/toc/03 derivatives/04 product and quotient rules \leaders\hbox to 10pt {\hss.\hss}  \hfill \xref{r58}
\smallskip
ink/07C \leaders\hbox to 10pt {\hss.\hss}  \hfill \xref{r59}
\smallskip
ink/07D \leaders\hbox to 10pt {\hss.\hss}  \hfill \xref{r60}
\smallskip
ink/07I \leaders\hbox to 10pt {\hss.\hss}  \hfill \xref{r61}
\smallskip
ink/07M \leaders\hbox to 10pt {\hss.\hss}  \hfill \xref{r62}
\smallskip
ink/07O \leaders\hbox to 10pt {\hss.\hss}  \hfill \xref{r63}
\smallskip
ink/07V \leaders\hbox to 10pt {\hss.\hss}  \hfill \xref{r64}
\smallskip
ink/07W \leaders\hbox to 10pt {\hss.\hss}  \hfill \xref{r65}
\smallskip
ink/07X \leaders\hbox to 10pt {\hss.\hss}  \hfill \xref{r66}
\smallskip
ink/07Y \leaders\hbox to 10pt {\hss.\hss}  \hfill \xref{r67}
\smallskip
ink/080 \leaders\hbox to 10pt {\hss.\hss}  \hfill \xref{r68}
\smallskip
\vfill \break
\xrdef{r1}
2 4 01
\smallskip
\par\begingroup
\leftskip=2em
path: calculus/pearson/examples/2 4 01
\smallskip
parents:  04 infinite limits [\xref{r51}]
\smallskip
\leftskip=0em
\par\endgroup
\bigskip
\xrdef{r2}
2 4 02
\smallskip
\par\begingroup
\leftskip=2em
path: calculus/pearson/examples/2 4 02
\smallskip
parents:  04 infinite limits [\xref{r51}]
\smallskip
\leftskip=0em
\par\endgroup
\bigskip
\xrdef{r3}
2 4 03
\smallskip
\par\begingroup
\leftskip=2em
path: calculus/pearson/examples/2 4 03
\smallskip
parents:  04 infinite limits [\xref{r51}]
\smallskip
\leftskip=0em
\par\endgroup
\bigskip
\xrdef{r4}
2 4 04
\smallskip
\par\begingroup
\leftskip=2em
path: calculus/pearson/examples/2 4 04
\smallskip
parents:  04 infinite limits [\xref{r51}]
\smallskip
\leftskip=0em
\par\endgroup
\bigskip
\xrdef{r5}
2 4 05
\smallskip
\par\begingroup
\leftskip=2em
path: calculus/pearson/examples/2 4 05
\smallskip
parents:  04 infinite limits [\xref{r51}]
\smallskip
\leftskip=0em
\par\endgroup
\bigskip
\xrdef{r6}
2 4 06
\smallskip
\par\begingroup
\leftskip=2em
path: calculus/pearson/examples/2 4 06
\smallskip
parents:  04 infinite limits [\xref{r51}]
\smallskip
\leftskip=0em
\par\endgroup
\bigskip
\xrdef{r7}
2 5 01
\smallskip
\par\begingroup
\leftskip=2em
path: calculus/pearson/examples/2 5 01
\smallskip
parents:  05 limits at infinity [\xref{r52}]
\smallskip
\leftskip=0em
\par\endgroup
\bigskip
\xrdef{r8}
2 5 02
\smallskip
\par\begingroup
\leftskip=2em
path: calculus/pearson/examples/2 5 02
\smallskip
parents:  05 limits at infinity [\xref{r52}]
\smallskip
\leftskip=0em
\par\endgroup
\bigskip
\xrdef{r9}
2 5 03
\smallskip
\par\begingroup
\leftskip=2em
path: calculus/pearson/examples/2 5 03
\smallskip
parents:  05 limits at infinity [\xref{r52}]
\smallskip
\leftskip=0em
\par\endgroup
\bigskip
\xrdef{r10}
2 5 04
\smallskip
\par\begingroup
\leftskip=2em
path: calculus/pearson/examples/2 5 04
\smallskip
parents:  05 limits at infinity [\xref{r52}]
\smallskip
\leftskip=0em
\par\endgroup
\bigskip
\xrdef{r11}
2 5 05
\smallskip
\par\begingroup
\leftskip=2em
path: calculus/pearson/examples/2 5 05
\smallskip
parents:  05 limits at infinity [\xref{r52}]
\smallskip
\leftskip=0em
\par\endgroup
\bigskip
\xrdef{r12}
2 5 06
\smallskip
\par\begingroup
\leftskip=2em
path: calculus/pearson/examples/2 5 06
\smallskip
parents:  05 limits at infinity [\xref{r52}]
\smallskip
\leftskip=0em
\par\endgroup
\bigskip
\xrdef{r13}
ink
\smallskip
\par\begingroup
\leftskip=2em
path: calculus/pearson/ink
\smallskip
children: 07V [\xref{r64}], 07W [\xref{r65}], 07Y [\xref{r67}], 080 [\xref{r68}], 07X [\xref{r66}], 07C [\xref{r59}], 07O [\xref{r63}], 07I [\xref{r61}], 07D [\xref{r60}], 07M [\xref{r62}]
\smallskip
\leftskip=0em
\par\endgroup
\bigskip
\xrdef{r14}
news
\smallskip
\par\begingroup
\leftskip=2em
path: calculus/pearson/news
\smallskip
\leftskip=0em
\par\endgroup
\bigskip
\xrdef{r15}
theorems
\smallskip
\par\begingroup
\leftskip=2em
path: calculus/pearson/theorems
\smallskip
\leftskip=0em
\par\endgroup
\bigskip
\xrdef{r16}
2 10 polynomial and rational functions
\smallskip
\par\begingroup
\leftskip=2em
path: calculus/pearson/theorems/2 10 polynomial and rational functions
\smallskip
text:  Polynomial and Rational Functions. \par {\bf a.} A polynomial function is continuous for all x.  \par {\bf b.} A rational function (a function of the form ${p \over q}$, where $p$ and $q$ are polynomials) is continuous for all $x$ for which $q(x) \ne 0$. 
\smallskip
\leftskip=0em
\par\endgroup
\bigskip
\xrdef{r17}
2 11 continunity of composite functions
\smallskip
\par\begingroup
\leftskip=2em
path: calculus/pearson/theorems/2 11 continunity of composite functions
\smallskip
text:  Continunity of Composite Functions at a Point. if $g$ is continuous at $a$ and $f$ is continuous at $g(a)$, then the composite function $g \circ f$ is continuous at $a$.
\smallskip
\leftskip=0em
\par\endgroup
\bigskip
\xrdef{r18}
2 12 limits of composite functions
\smallskip
\par\begingroup
\leftskip=2em
path: calculus/pearson/theorems/2 12 limits of composite functions
\smallskip
text:  Limits of Composite Functions. \par {\bf 1.} if $g$ is continuous at $a$ and $f$ is continuous at $g(a)$, then  $$\ \lim_{x \to a} f(g(x)) = f\big(\lim_{x \to a} g(x)\big). $$  \par {\bf 2.} if $g(x) = L$ and $f$ is continuous at $L$, then  $$\ \lim_{x \to a} f(g(x)) = f\big(\lim_{x \to a} g(x)\big). $$ 
\smallskip
\leftskip=0em
\par\endgroup
\bigskip
\xrdef{r19}
2 13 continuity of functions with roots
\smallskip
\par\begingroup
\leftskip=2em
path: calculus/pearson/theorems/2 13 continuity of functions with roots
\smallskip
text:  Continuity of Functions with Roots. Assume $n$ is a positive integer. If $n$ is an odd integer, then $(f(x))^{1/n}$ is continuous at all points at which $f$ is continuous. \par If $n$ is even, then $(f(x))^{1/n}$ is continuous at all points $a$ at which $f$ is continuous and $f(a) > 0$.
\smallskip
\leftskip=0em
\par\endgroup
\bigskip
\xrdef{r20}
2 14 continuity of inverse functions
\smallskip
\par\begingroup
\leftskip=2em
path: calculus/pearson/theorems/2 14 continuity of inverse functions
\smallskip
text:  Continuity of Inverse Functions. If a function $f$ is continuous on an interval $I$ and has an inverse on $I$, then its inverse $f^{-1}$ is also continuous (on the interval consisting of the points $f(x)$, where $x$ is in $I$).
\smallskip
\leftskip=0em
\par\endgroup
\bigskip
\xrdef{r21}
2 15 continuity of transcendental
\smallskip
\par\begingroup
\leftskip=2em
path: calculus/pearson/theorems/2 15 continuity of transcendental
\smallskip
text:  Continuity of Transcendental functions. TODO.
\smallskip
\leftskip=0em
\par\endgroup
\bigskip
\xrdef{r22}
2 16 intermediate value theorem
\smallskip
\par\begingroup
\leftskip=2em
path: calculus/pearson/theorems/2 16 intermediate value theorem
\smallskip
text:  Intermediate Value Theorem. Suppose $f$ is continuous on the interval $[a, b]$ and $L$ is a number strictly between $f(a)$ and $f(b)$. Then there exists at lest one number $c$ in $(a, b)$ satisfying $f(c) = L$.
\smallskip
\leftskip=0em
\par\endgroup
\bigskip
\xrdef{r23}
2 1 relationship between one-sided
\smallskip
\par\begingroup
\leftskip=2em
path: calculus/pearson/theorems/2 1 relationship between one-sided
\smallskip
text:  Relationship between One-Sided and Two-Sided Limits. Assume $f$ is defined for all $x$ near $a$ except possibly at $a$. Then  $$\ \lim_{x \to a} f(x) = L $$  if and only if  $$\ \lim_{x \to a^{-}} = L $$  and  $$\ \lim_{x \to a^{+}} = L $$
\smallskip
\leftskip=0em
\par\endgroup
\bigskip
\xrdef{r24}
2 2 limits of a linear function
\smallskip
\par\begingroup
\leftskip=2em
path: calculus/pearson/theorems/2 2 limits of a linear function
\smallskip
text:  Limits of a Linear Function. Let $a$, $b$, and $m$ be real numbers. For linear functions $f(x) = mx + b$,  $$\ \lim_{x \to a} {f(x)} = f(a) = ma + b. $$
\smallskip
\leftskip=0em
\par\endgroup
\bigskip
\xrdef{r25}
2 3a limit law: sum
\smallskip
\par\begingroup
\leftskip=2em
path: calculus/pearson/theorems/2 3a limit law: sum
\smallskip
text:  Limit Law: Sum. $$\ \lim_{x \to a}{ (f(x) + g(x)) } = \lim_{x \to a}{f(x)} + \lim_{x \to a}{g(x)}  $$
\smallskip
\leftskip=0em
\par\endgroup
\bigskip
\xrdef{r26}
2 3b limit law: difference
\smallskip
\par\begingroup
\leftskip=2em
path: calculus/pearson/theorems/2 3b limit law: difference
\smallskip
text:  Limit Law: Difference. $$\ \lim_{x \to a}{ (f(x) - g(x)) } = \lim_{x \to a}{f(x)} - \lim_{x \to a}{g(x)}  $$
\smallskip
\leftskip=0em
\par\endgroup
\bigskip
\xrdef{r27}
2 3c limit law: constant multiple
\smallskip
\par\begingroup
\leftskip=2em
path: calculus/pearson/theorems/2 3c limit law: constant multiple
\smallskip
text:  Limit Law: Constant multiple. $$\ \lim_{x \to a}{ cf(x) } = c\lim_{x \to a}{f(x)} $$
\smallskip
\leftskip=0em
\par\endgroup
\bigskip
\xrdef{r28}
2 3d limit law: product
\smallskip
\par\begingroup
\leftskip=2em
path: calculus/pearson/theorems/2 3d limit law: product
\smallskip
text:  Limit Law: Product. $$\ \lim_{x \to a}{ (f(x)g(x)) } = \bigg( \lim_{x \to a}{f(x)} \bigg) \bigg( \lim_{x \to a}{g(x)}  \bigg) $$
\smallskip
\leftskip=0em
\par\endgroup
\bigskip
\xrdef{r29}
2 3e limit law: quotient
\smallskip
\par\begingroup
\leftskip=2em
path: calculus/pearson/theorems/2 3e limit law: quotient
\smallskip
text:  Limit Law: Quotient. $$\ \lim_{x \to a}{ \bigg( { f(x) \over g(x) } \bigg) } = { {\lim_{x \to a}{f(x)}} \over {\lim_{x \to a}{g(x)}} } $$  provided  $$\ \lim_{x \to a}{g(x)} \ne 0 $$
\smallskip
\leftskip=0em
\par\endgroup
\bigskip
\xrdef{r30}
2 3f limit law: power
\smallskip
\par\begingroup
\leftskip=2em
path: calculus/pearson/theorems/2 3f limit law: power
\smallskip
text:  Limit Law: Power. $$\ \lim_{x \to a}{ (f(x))^n } = \big( \lim_{x \to a}{ (f(x)) } \big)^n $$
\smallskip
\leftskip=0em
\par\endgroup
\bigskip
\xrdef{r31}
2 3g limit law: root
\smallskip
\par\begingroup
\leftskip=2em
path: calculus/pearson/theorems/2 3g limit law: root
\smallskip
text:  Limit Law: Root. $$\ \lim_{x \to a}{ (f(x))^{1/n} } = \big( \lim_{x \to a}{ (f(x)) } \big)^{1/n} $$  provided $f(x) > 0$ for $x$ near $a$, if $n$ is even.
\smallskip
\leftskip=0em
\par\endgroup
\bigskip
\xrdef{r32}
2 3h limit laws for one-sided limits
\smallskip
\par\begingroup
\leftskip=2em
path: calculus/pearson/theorems/2 3h limit laws for one-sided limits
\smallskip
text:  Limit laws for One-Sided Limits. Laws 2.3a-2.3f hold with $\lim\limits_{x \to a}$ replaced with $\lim\limits_{x \to a^+}$ or $\lim\limits_{x \to a^-}$. Law 2.3g is modified as follows.  \par a.  $$\ \lim_{x \to a^+}(f(x))^{1/n} = \bigg( \lim_{x \to a^+} f(x) \bigg)^{1/n} $$  provided $f(x) \ge 0$, for $x$ near $a$ with $x > a$, if $n$ is even.  \par b.  $$\ \lim_{x \to a^-}(f(x))^{1/n} = \bigg( \lim_{x \to a^-} f(x) \bigg)^{1/n} $$  provided $f(x) \ge 0$, for $x$ near $a$ with $x < a$, if $n$ is even. 
\smallskip
\leftskip=0em
\par\endgroup
\bigskip
\xrdef{r33}
2 4 limits of polynomial and rational
\smallskip
\par\begingroup
\leftskip=2em
path: calculus/pearson/theorems/2 4 limits of polynomial and rational
\smallskip
text:  Limits of Polynomial and Rational Functions. Assume $p$ and $q$ are polynomials and $a$ is a constant.  \par a. Polynomial functions $$\ \lim_{x \to a} p(x) = p(a) $$   \par b. Rational functions  $$\ \lim_{x \to a} {p(x) \over q(x)} = {p(a) \over q(a)} $$  provided $q(a) \ne 0$.
\smallskip
\leftskip=0em
\par\endgroup
\bigskip
\xrdef{r34}
2 5 the squeeze theorem
\smallskip
\par\begingroup
\leftskip=2em
path: calculus/pearson/theorems/2 5 the squeeze theorem
\smallskip
text:  The Squeeze Theorem. Assume the functions $f$, $g$, and $h$ satisfy $f(x) \le g(x) \le h(x)$ for all values of $x$ near $a$, except possibly at $a$. If $\lim\limits_{x \to a} f(x) = \lim\limits_{x \to a} h(x) = L$, then $\lim\limits_{x \to a} g(x) = L$.
\smallskip
\leftskip=0em
\par\endgroup
\bigskip
\xrdef{r35}
2 6 limits at infinity of powers and
\smallskip
\par\begingroup
\leftskip=2em
path: calculus/pearson/theorems/2 6 limits at infinity of powers and
\smallskip
text:  Limits at Infinity of Powers and Polynomials. Let $n$ be a positive integer and let $p$ be the polynomial  $p(x) = a_n x^n + a_{n - 1}x^{n - 1} + \cdots + a_2 x^2 + a_1x + a_0$, where $a_n \ne 0$. \par 1. $\lim\limits_{x \to \pm\infty} x^n = \infty$ when $n$ is even.  \par 2. $\lim\limits_{x \to \infty} x^n = \infty$ and $\lim\limits_{x \to -\infty} x^n = -\infty$ when $n$ is odd.  \par 3. $\lim\limits_{x \to \pm\infty} {1 \over x^n} = \lim\limits_{x \to \pm\infty} {x^{-n}} = 0 $.  \par 4. $\lim\limits_{x \to \pm\infty} p(x)= \lim\limits_{x \to \pm\infty} a_nx^n= \pm\infty$, depending on the degree of the polynomial and the sign of the leading coefficient $a_n$. 
\smallskip
\leftskip=0em
\par\endgroup
\bigskip
\xrdef{r36}
2 7 end behavior and asymptotes of
\smallskip
\par\begingroup
\leftskip=2em
path: calculus/pearson/theorems/2 7 end behavior and asymptotes of
\smallskip
text:  End Behavior and Asymptotes of Rational Functions. Suppose $f(x) = {p(x) \over q(x)}$ is a rational function, where $$\ p(x) = a_mx^m + a_{m-1}x^{m-1} + \cdots + a_2x^2 + a_0 $$  and  $$\ q(x) = b_nx^n + b_{n-1}x^{n-1} + \cdots + b_2x^2 + b_0 $$  with $a_m \ne 0$ and $b_n \ne 0$.  \par {\bf a. Degree of numerator less than degree of denominator.} If $m < n$, then $\lim\limits_{x \to \pm\infty} f(x) = 0$, and $y = 0$ is a horizontal asymptote of $f$.  \par {\bf b. Degree of numerator equals degree of denominator.} If $m = n$, then $\lim\limits_{x \to \pm\infty} f(x) = a_m/b_n$, and $y = a_m/b_n$ is a horizontal asymptote of $f$.  \par {\bf c. Degree of numerator greater than degree of denominator.} If $m > n$, then $\lim\limits_{x \to \pm\infty} f(x) = \infty$ or $-\infty$, and $f$ has no horizontal asymptote.  \par {\bf d. Slant asymptote.}  If $m = n + 1$, then $\lim\limits_{x \to \pm\infty} f(x) = \infty$ or $-\infty$, and $f$ has no horizontal asymptote, but $f$ has a slant asymptote.  \par {\bf e. Vertifcal asymptotes.} Assuming $f$ is in reduced form ($p$ and $q$ share no common factors), vertical asymptotes occur at the zeros of $q$.
\smallskip
\leftskip=0em
\par\endgroup
\bigskip
\xrdef{r37}
2 8 end behavior of e x, e {-x}, and
\smallskip
\par\begingroup
\leftskip=2em
path: calculus/pearson/theorems/2 8 end behavior of e x, e {-x}, and
\smallskip
text:  End Behavior of $e^x$, $e^{-x}$, and $\ln x$. The end behavior for $e^x$ and $e^{-x}$ on $(-\infty, infty)$ and on $\ln x$ on $(0, \infty)$ is given by the following limits:  $$\ \lim_{x \to \infty} e^x = \infty, \lim_{x \to -\infty} e^x = 0 $$  $$\ \lim_{x \to \infty} e^{-x} = 0, \lim_{x \to -\infty} e^{-x} = \infty $$  $$\ \lim_{x \to 0^+} \ln x = -\infty, \lim_{x \to \infty} \ln x = \infty $$
\smallskip
\leftskip=0em
\par\endgroup
\bigskip
\xrdef{r38}
2 9 continuity rules
\smallskip
\par\begingroup
\leftskip=2em
path: calculus/pearson/theorems/2 9 continuity rules
\smallskip
text:  Continuity Rules. If $f$ and $g$ are continuous at $a$, then the following functions are also continuous at $a$. Assume $c$ is a constant and $n > 0$ is an integer.  \par {\bf a.} $f + g$  \par {\bf b.} $f - g$  \par {\bf c.} $cf$  \par {\bf d.} $fg$  \par {\bf e.} $f/g$, provided $g(a) \ne 0$  \par {\bf f.} $(f(x))^n$ 
\smallskip
\leftskip=0em
\par\endgroup
\bigskip
\xrdef{r39}
3 1 differentiable implies continuous
\smallskip
\par\begingroup
\leftskip=2em
path: calculus/pearson/theorems/3 1 differentiable implies continuous
\smallskip
text:  Differentiable Implies Continuous. If $f$ is differentiable at $a$, then $f$ is continuous at $a$.
\smallskip
\leftskip=0em
\par\endgroup
\bigskip
\xrdef{r40}
3 1a not continuous implies not differentiable
\smallskip
\par\begingroup
\leftskip=2em
path: calculus/pearson/theorems/3 1a not continuous implies not differentiable
\smallskip
text:  Not Continuous implies Not Differentiable. If $f$ is not continuous at $a$, then $f$ is not differentiable at $a$.
\smallskip
\leftskip=0em
\par\endgroup
\bigskip
\xrdef{r41}
3 2 constant rule
\smallskip
\par\begingroup
\leftskip=2em
path: calculus/pearson/theorems/3 2 constant rule
\smallskip
text:  Constant Rule. If $c$ is a real number, then ${d \over dx}(c) = 0$.
\smallskip
\leftskip=0em
\par\endgroup
\bigskip
\xrdef{r42}
3 3 power rule
\smallskip
\par\begingroup
\leftskip=2em
path: calculus/pearson/theorems/3 3 power rule
\smallskip
text:  Power Rule. If $n$ is a nonnegative integer, then ${d \over dx}(x^n) = nx^{n - 1}$
\smallskip
\leftskip=0em
\par\endgroup
\bigskip
\xrdef{r43}
3 4 constant multiple rule
\smallskip
\par\begingroup
\leftskip=2em
path: calculus/pearson/theorems/3 4 constant multiple rule
\smallskip
text:  Constant Multiple Rule. If $f$ is differentiable at $x$ and $c$ is a constant, then $$\ {d \over dx}(cf(x)) = c f'(x). $$
\smallskip
\leftskip=0em
\par\endgroup
\bigskip
\xrdef{r44}
3 5 sum rule
\smallskip
\par\begingroup
\leftskip=2em
path: calculus/pearson/theorems/3 5 sum rule
\smallskip
text:  Sum Rule. If $f$ and $g$ are differentiable at $x$, then  $$\ {d \over dx}(f(x) + g(x)) = f'(x) + g'(x). $$
\smallskip
\leftskip=0em
\par\endgroup
\bigskip
\xrdef{r45}
3 6 the derivative of e x
\smallskip
\par\begingroup
\leftskip=2em
path: calculus/pearson/theorems/3 6 the derivative of e x
\smallskip
text:  The Derivative of $e^x$. The function $f(x) = e^x$ is differentiable for all real numbers $x$, and $$\ {d \over dx}(e^x) = e^x $$
\smallskip
\leftskip=0em
\par\endgroup
\bigskip
\xrdef{r46}
3 7 product rule
\smallskip
\par\begingroup
\leftskip=2em
path: calculus/pearson/theorems/3 7 product rule
\smallskip
text:  Product Rule. If $f$ and $g$ are differentiable at $x$, then  $$\ {d \over dx}(f(x)g(x)) = f'(x)g(x) + f(x)g'(x). $$
\smallskip
\leftskip=0em
\par\endgroup
\bigskip
\xrdef{r47}
3 8 quotient rule
\smallskip
\par\begingroup
\leftskip=2em
path: calculus/pearson/theorems/3 8 quotient rule
\smallskip
text:  Quotient Rule. If $f$ and $g$ are differntiable at $x$ and $g(x) \ne 0$, then the derivative of $f / g$ at $x$ exists and $$\ {d \over dx}\bigg({f(x) \over g(x)}\bigg) = { g(x)f'(x) - f(x)g'(x) \over (g(x))^2 } $$
\smallskip
\leftskip=0em
\par\endgroup
\bigskip
\xrdef{r48}
3 9 power rule
\smallskip
\par\begingroup
\leftskip=2em
path: calculus/pearson/theorems/3 9 power rule
\smallskip
text:  Power Rule (general form). If $n$ is any real number, then $$\ {d \over dx}(x^n) = nx^{n - 1} $$
\smallskip
\leftskip=0em
\par\endgroup
\bigskip
\xrdef{r49}
toc
\smallskip
\par\begingroup
\leftskip=2em
path: calculus/pearson/toc
\smallskip
\leftskip=0em
\par\endgroup
\bigskip
\xrdef{r50}
02 limits
\smallskip
\par\begingroup
\leftskip=2em
path: calculus/pearson/toc/02 limits
\smallskip
\leftskip=0em
\par\endgroup
\bigskip
\xrdef{r51}
04 infinite limits
\smallskip
\par\begingroup
\leftskip=2em
path: calculus/pearson/toc/02 limits/04 infinite limits
\smallskip
children: 2 4 01 [\xref{r1}], 2 4 03 [\xref{r3}], 2 4 02 [\xref{r2}], 2 4 05 [\xref{r5}], 2 4 06 [\xref{r6}], 2 4 04 [\xref{r4}]
\smallskip
\leftskip=0em
\par\endgroup
\bigskip
\xrdef{r52}
05 limits at infinity
\smallskip
\par\begingroup
\leftskip=2em
path: calculus/pearson/toc/02 limits/05 limits at infinity
\smallskip
children: 2 5 04 [\xref{r10}], 2 5 06 [\xref{r12}], 2 5 05 [\xref{r11}], 2 5 02 [\xref{r8}], 2 5 03 [\xref{r9}], 2 5 01 [\xref{r7}]
\smallskip
\leftskip=0em
\par\endgroup
\bigskip
\xrdef{r53}
06 continuity
\smallskip
\par\begingroup
\leftskip=2em
path: calculus/pearson/toc/02 limits/06 continuity
\smallskip
\leftskip=0em
\par\endgroup
\bigskip
\xrdef{r54}
03 derivatives
\smallskip
\par\begingroup
\leftskip=2em
path: calculus/pearson/toc/03 derivatives
\smallskip
\leftskip=0em
\par\endgroup
\bigskip
\xrdef{r55}
01 introducing the derivative
\smallskip
\par\begingroup
\leftskip=2em
path: calculus/pearson/toc/03 derivatives/01 introducing the derivative
\smallskip
\leftskip=0em
\par\endgroup
\bigskip
\xrdef{r56}
02 derivative as function
\smallskip
\par\begingroup
\leftskip=2em
path: calculus/pearson/toc/03 derivatives/02 derivative as function
\smallskip
\leftskip=0em
\par\endgroup
\bigskip
\xrdef{r57}
03 rules of differentiation
\smallskip
\par\begingroup
\leftskip=2em
path: calculus/pearson/toc/03 derivatives/03 rules of differentiation
\smallskip
\leftskip=0em
\par\endgroup
\bigskip
\xrdef{r58}
04 product and quotient rules
\smallskip
\par\begingroup
\leftskip=2em
path: calculus/pearson/toc/03 derivatives/04 product and quotient rules
\smallskip
\leftskip=0em
\par\endgroup
\bigskip
\xrdef{r59}
07C
\smallskip
\par\begingroup
\leftskip=2em
path: ink/07C
\smallskip
parents:  ink [\xref{r13}]
\smallskip
\leftskip=0em
\par\endgroup
\bigskip
\xrdef{r60}
07D
\smallskip
\par\begingroup
\leftskip=2em
path: ink/07D
\smallskip
parents:  ink [\xref{r13}]
\smallskip
\leftskip=0em
\par\endgroup
\bigskip
\xrdef{r61}
07I
\smallskip
\par\begingroup
\leftskip=2em
path: ink/07I
\smallskip
parents:  ink [\xref{r13}]
\smallskip
\leftskip=0em
\par\endgroup
\bigskip
\xrdef{r62}
07M
\smallskip
\par\begingroup
\leftskip=2em
path: ink/07M
\smallskip
parents:  ink [\xref{r13}]
\smallskip
\leftskip=0em
\par\endgroup
\bigskip
\xrdef{r63}
07O
\smallskip
\par\begingroup
\leftskip=2em
path: ink/07O
\smallskip
parents:  ink [\xref{r13}]
\smallskip
\leftskip=0em
\par\endgroup
\bigskip
\xrdef{r64}
07V
\smallskip
\par\begingroup
\leftskip=2em
path: ink/07V
\smallskip
parents:  ink [\xref{r13}]
\smallskip
\leftskip=0em
\par\endgroup
\bigskip
\xrdef{r65}
07W
\smallskip
\par\begingroup
\leftskip=2em
path: ink/07W
\smallskip
parents:  ink [\xref{r13}]
\smallskip
\leftskip=0em
\par\endgroup
\bigskip
\xrdef{r66}
07X
\smallskip
\par\begingroup
\leftskip=2em
path: ink/07X
\smallskip
parents:  ink [\xref{r13}]
\smallskip
\leftskip=0em
\par\endgroup
\bigskip
\xrdef{r67}
07Y
\smallskip
\par\begingroup
\leftskip=2em
path: ink/07Y
\smallskip
parents:  ink [\xref{r13}]
\smallskip
\leftskip=0em
\par\endgroup
\bigskip
\xrdef{r68}
080
\smallskip
\par\begingroup
\leftskip=2em
path: ink/080
\smallskip
parents:  ink [\xref{r13}]
\smallskip
\leftskip=0em
\par\endgroup
\bigskip
\bye
